% Noether P09 Korean producer draft — UNCHECKED
% T02-U13; ED0002 lines6495--6511; source 1,266 LF B / 2D046108BEE05CA9E3C0A4026DDE1FC616FB1FA44577AC489E36DA1481FDC7A4
% Pointer v009 / ED0002: B06BE3530D9CF2E82B56FDBA7FE41D5D044DF2425DFA2A059D4939EAA2F7A6C2 / C9A125167ACB33D914EE4374B65AE7CDF0052F568371B8B77B720EA178ABF0E3
% Translation only; Teiler is provisionally rendered by the neutral relation 부분정역으로 포함한다.

$\mathfrak G$를 추상적 성질 I--IV를 만족하는, 정수적 초월수로 이루어진 임의의 주어진 영역이라 하자. E. 체르멜로가 모든 복소수의 체에 적용한 절차에 따라 $\mathfrak G$에서 \emph{$\mathfrak G$의 모든 수의 초월기저} $H$를 골라낸다. II에 따르면 \emph{모든} 실수 또는 복소수를 $\mathfrak G$의 두 수의 몫으로 나타낼 수 있으므로, $H$는 동시에 \emph{모든} 수의 초월기저이다. 따라서 주어진 모든 영역 $\mathfrak G$에 대하여 모든 수의 초월기저를 $\mathfrak G$에 속하도록 선택할 수 있고, 그 결과 \emph{기저 성질 1)이 성립한다.} 또한 III에 따라 $\mathfrak G$는 모든 대수적 정수의 정역 $[K]$를 포함하며, 따라서 I에 따라 기저의 수 $\eta$에 관한, 계수가 $[K]$에 속하는 모든 다항식, 즉 $\eta$의 모든 정수적 다항식도 포함한다. 이 다항식들은 정역 $[H]$를 이룬다. 그러므로 \emph{기저 성질 2)도 언제나 성립한다}. 다시 말해, \emph{모든 영역 $\mathfrak G$는 모든 수의 초월기저 $H$를 이용하여, $\mathfrak G$가 정역 $[H]$를 부분정역으로 포함하도록 구성할 수 있다.}
